\clearpage
\setcounter{page}{1}
\maketitlesupplementary

\section{Paired-Intervention Ablation}
\label{sec:supp_paired}

This supplement expands the controlled comparison in
\cref{sec:paired_ablation}. All results concern the 10-category core
dataset and frozen Qwen2.5-VL-3B-Instruct. The larger-menu
paired-subtraction experiment has not been completed.

\subsection{Endpoint Contributions and Fitting}

For a record containing categories \(c\) and \(q\), let
\(H_b,H_c,H_q\in\mathbb R^{T\times D}\) be the aligned visual-token
states of the base and two composites. Shared and endpoint-specific
weights are
\begin{equation}
    \begin{aligned}
        e_t&=\tfrac12\bigl(\|H_{c,t}-H_{b,t}\|_2
                         +\|H_{q,t}-H_{b,t}\|_2\bigr),\\
        e_{z,t}&=\tfrac12\|H_{z,t}-H_{b,t}\|_2,
        \quad z\in\{c,q\},\\
        \alpha_t&=\frac{e_t}{\sum_s e_s+10^{-8}},\qquad
        \alpha_{z,t}=\frac{e_{z,t}}{\sum_s e_{z,s}+10^{-8}}.
    \end{aligned}
\end{equation}
\Cref{tab:paired_endpoints} specifies both endpoint contributions.
An endpoint-specific weight uses only its own composite and the base;
the half factor and numerical stabilizer are retained in all variants.

\begin{table*}[t]
    \centering
    \small
    \begin{tabular}{lll}
        \toprule
        Estimator & Contribution to category \(c\) & Contribution to category \(q\) \\
        \midrule
        P0: paired, shared weights
            & \(\alpha^\top(H_c-H_q)\) & \(\alpha^\top(H_q-H_c)\) \\
        P1: single-sided, shared weights
            & \(\alpha^\top(H_c-H_b)\) & \(\alpha^\top(H_q-H_b)\) \\
        P2: raw, shared weights
            & \(\alpha^\top H_c\) & \(\alpha^\top H_q\) \\
        P3: single-sided, own weights
            & \(\alpha_c^\top(H_c-H_b)\) & \(\alpha_q^\top(H_q-H_b)\) \\
        \bottomrule
    \end{tabular}
    \caption{\textbf{Exact endpoint contributions.} All estimators use the
        same records. Only P0 is antisymmetric between endpoints. P2 retains
        intervention-derived weights despite omitting subtraction from
        the pooled representation.}
    \label{tab:paired_endpoints}
\end{table*}

Let \(\mathcal R_c\) contain every record in which \(c\) appears at
either endpoint, and let \(x_{i,c}^{(v)}\) be the corresponding
contribution for variant \(v\). The fitted direction is
\begin{equation}
    \widetilde g_c^{(v)}=\frac{1}{|\mathcal R_c|}
        \sum_{i\in\mathcal R_c}x_{i,c}^{(v)},\qquad
    g_c^{(v)}=\frac{\widetilde g_c^{(v)}}{
        \|\widetilde g_c^{(v)}\|_2}.
\end{equation}
We introduce no per-record normalization, centering, or
competitor-frequency correction. The 2,104 records produce 4,208 endpoint
contributions across 45 observed pairs; pair counts range from 4 to 101.
\Cref{tab:paired_classes} lists the per-category fitting and evaluation
counts. Consequently, this ablation tests the estimator on its observed
competitor frequencies rather than an exactly balanced menu.

\subsection{Data Replay and Numerical Settings}

All variants use the same 1,324 referenced donor cutouts and 2,104 saved
recipes, with COCO train backgrounds and COCO val2017 evaluation images.
We recover the saved donor identities, paste sites, and area parameters
without resampling donors or rerunning selection gates. Compositing uses
3-pixel feathering and a 16-pixel dilation for local color harmonization.
The resizing routine matches the target \emph{bounding-box area} while
preserving aspect ratio; realized foreground-mask areas can differ between
the two cutouts. Placement uses a shared requested center with boundary
clipping. Matching these construction parameters does not guarantee
identical edit geometry or context effects.

We use model revision
\texttt{66285546d2b8}
with the fixed prompt \texttt{Describe.}, BF16 inference, and FP16
hidden-state storage. Contributions and cosine scores use FP32 arithmetic;
class contributions are accumulated in FP64 before normalization.
The response-weight stabilizer is \(10^{-8}\), while the cosine norm
floor is \(10^{-12}\). The implementation rejects empty categories and
zero or nonfinite class means. No extraction records failed.
Hidden states are taken at layers 1, 9, 18, 27, and 35 using the same
extraction convention for discovery and evaluation, including final
normalization at the last layer. Extraction uses two H100 GPUs.

\subsection{Uncertainty and Random Controls}

For each layer, we resample 5,000 evaluation-image indices with replacement
for 2,000 bootstrap draws. Each draw is shared across all categories and
estimators, preserving within-image dependence. We recompute each
category's AUC, average over the fixed 10-category menu, and form the
paired difference within each draw. Intervals are the 2.5th and 97.5th
percentiles. All draws are valid for every reported comparison. The
bootstrap seed is \(70000+\ell\).

The random control uses 25 independent sets of unit Gaussian directions
per layer, with seeds \(3000+100\ell+s\), \(s=0,\ldots,24\).
We average their metrics, not their scores. At L18, the mean random macro
AUC is 0.5203 and its standard deviation across sets is 0.0365.
For comparisons to this control, bootstrap draws retain the same 25
direction sets and average their resampled metrics. This seed standard
deviation is distinct from the image-bootstrap interval.

The primary endpoint is P0 minus P1 macro AUC at L18, with a practical
difference threshold of 0.005 AUC. The protocol was specified before this
run but after inspection of earlier validation results. These are
follow-up analyses on an existing evaluation population. Intervals
condition on the fixed fitting records and directions; they do not
incorporate new donor, background, or gate realizations. AP is reported
as a point estimate. Other layer and category comparisons are diagnostic;
their intervals have no multiplicity correction.

% Generated by scripts/render_paired_ablation.py; do not edit numbers.
\begin{table*}[t]
    \centering
    \footnotesize
    \setlength{\tabcolsep}{4pt}
    \begin{tabular}{rrrrrrr}
        \toprule
        Layer & P0 & P1 & P2 & P3 & Random & P0$-$P1 (pp) [95\% CI] \\
        \midrule
        1 & 0.9543 & 0.8797 & 0.5778 & 0.8766 & 0.5133 & $+7.46\;[+6.72,+8.15]$ \\
        9 & 0.9622 & 0.8895 & 0.5658 & 0.8858 & 0.5175 & $+7.27\;[+6.61,+7.96]$ \\
        18 & 0.9606 & 0.8818 & 0.5102 & 0.8804 & 0.5203 & $+7.88\;[+7.17,+8.60]$ \\
        27 & 0.9520 & 0.8443 & 0.5099 & 0.8439 & 0.5182 & $+10.77\;[+9.85,+11.72]$ \\
        35 & 0.9514 & 0.9009 & 0.5333 & 0.8980 & 0.4968 & $+5.05\;[+4.49,+5.60]$ \\
        \bottomrule
    \end{tabular}
    \caption{\textbf{Macro ROC-AUC across layers.} Settings are identical to \cref{tab:paired_primary}. AUCs use the 0--1 scale; paired differences and interval endpoints use percentage points (pp). L18 is primary. Other layers are secondary and intervals are unadjusted for multiple comparisons.}
    \label{tab:paired_layers}
\end{table*}

% Generated by scripts/render_paired_ablation.py; do not edit numbers.
\begin{table}[t]
    \centering
    \footnotesize
    \setlength{\tabcolsep}{4pt}
    \begin{tabular}{rrrrrr}
        \toprule
        Layer & P0 & P1 & P2 & P3 & Random \\
        \midrule
        1 & 0.7446 & 0.4514 & 0.1243 & 0.4467 & 0.1124 \\
        9 & 0.7939 & 0.4752 & 0.1150 & 0.4693 & 0.1126 \\
        18 & 0.7993 & 0.4504 & 0.1057 & 0.4494 & 0.1136 \\
        27 & 0.7653 & 0.3996 & 0.1020 & 0.4005 & 0.1157 \\
        35 & 0.7505 & 0.4315 & 0.1055 & 0.4248 & 0.1073 \\
        \bottomrule
    \end{tabular}
    \caption{\textbf{Macro average precision across layers.} Values are point estimates on the 0--1 scale. The random control averages AP over the same 25 direction sets used for AUC.}
    \label{tab:paired_ap}
\end{table}

% Generated by scripts/render_paired_ablation.py; do not edit numbers.
\begin{table}[t]
    \centering
    \footnotesize
    \setlength{\tabcolsep}{4pt}
    \begin{tabular}{lrrrrr}
        \toprule
        Category & Fits & Pos. & P0 & P1 & $\Delta$ (pp) \\
        \midrule
        car & 551 & 535 & 0.9789 & 0.9582 & $+2.07$ \\
        bus & 347 & 189 & 0.9930 & 0.9412 & $+5.18$ \\
        dog & 534 & 177 & 0.9840 & 0.8392 & $+14.48$ \\
        cat & 531 & 184 & 0.9901 & 0.9131 & $+7.70$ \\
        bird & 556 & 125 & 0.9608 & 0.8600 & $+10.08$ \\
        chair & 286 & 580 & 0.8396 & 0.7507 & $+8.89$ \\
        couch & 152 & 195 & 0.9648 & 0.8170 & $+14.78$ \\
        person & 587 & 2693 & 0.9669 & 0.8265 & $+14.04$ \\
        tv & 70 & 207 & 0.9862 & 0.9744 & $+1.18$ \\
        bottle & 594 & 379 & 0.9416 & 0.9375 & $+0.40$ \\
        \bottomrule
    \end{tabular}
    \caption{\textbf{Per-category AUC at L18.} Fits counts records containing each category at either endpoint; Pos. counts positive images among the 5,000 evaluation images. $\Delta$ is P0 minus P1 in percentage points. These are diagnostic point estimates; the primary endpoint is macro AUC.}
    \label{tab:paired_classes}
\end{table}


\subsection{Layer and Category Results}

P0 exceeds P1 in macro AUC and AP at all five layers
(\cref{tab:paired_layers}). The P3-minus-P1 AUC differences range from
\(-0.00365\) to \(-0.00035\); each corresponding 95\% interval
lies within the prespecified \(\pm0.005\) practical band. The raw
variant P2 has no clear advantage over the random control at L18. This
observation is layer-specific: its AUC exceeds the random-control mean
at L1, L9, and L35, with positive paired intervals.

At L18, P0 improves AUC over P1 for every category
(\cref{tab:paired_classes}). The largest point-estimate gains are for
couch (+14.78 points), dog (+14.48), and person (+14.04). These category
results do not resolve the separate chair/couch degradation observed with
a different, larger concept menu and different discovery records.

\subsection{Replay and Implementation Validation}

All 2,104 replayed records match the historical paste pixel counts, token
grids, and footprint-token counts. The recorded validation suites contain
146 passing tests, including endpoint-swap invariance and independence
of P3 weights from the opposite composite. Pilot extraction hooks agree
exactly with reference hidden states, with the appropriate final-layer
normalization. On identical tensors, P0 direction aggregation agrees
exactly with the original estimator at all five layers. Across all
5,000 evaluation images, the maximum per-category AUC discrepancy between
GPU scoring and the original CPU readout is \(5.71\times10^{-7}\),
below the \(10^{-5}\) acceptance threshold.

At L18, replayed and historical category directions have mean cosine
similarity 0.999914 and minimum 0.999829. Historical directions evaluated
on the same rebuilt image cache reach 0.960392 macro AUC, compared with
0.960574 for replayed P0, a difference of 0.000181. This establishes close
agreement of the averaged directions. It does not establish bitwise
recovery: record-level contrasts have median cosine 0.996900 and median
relative \(\ell_2\) difference 8.05\%, and the historical model revision
and donor-content hashes are unavailable.

The experiment therefore isolates the empirical effect of the subtraction
reference within a fixed, replayed pipeline. It neither identifies a
specific nuisance-cancellation mechanism nor tests a fully single-category
data-construction pipeline. Replication on the 76-class menu requires
recovering its original three-image caches or complete successful replay
recipes, including placement sites.
